\section{Campo magnético}

\subsection{Fuerza magnética sobre una carga en movimiento}
\begin{equation}
    \vec{F} = q \vec{v} \times \vec{B}
\end{equation}
donde \(q\) es la carga, \(\vec{v}\) su velocidad y \(\vec{B}\) el campo magnético.

\subsection{Magnitud y dirección}
\begin{equation}
    F = q v B \sin \theta
\end{equation}
La dirección está dada por la regla de la mano derecha para el producto cruz.

\section{Campo magnético producido por una corriente}

\subsection{Ley de Biot-Savart}
\begin{equation}
    d\vec{B} = \frac{\mu_0}{4\pi} \frac{I \, d\vec{l} \times \hat{r}}{r^2}
\end{equation}

\subsection{Campo de una espira circular}
\begin{equation}
    B = \frac{\mu_0 I R^2}{2 (R^2 + x^2)^{3/2}} \quad (\text{en el eje de la espira})
\end{equation}

\subsection{Campo de un solenoide ideal}
\begin{equation}
    B = \mu_0 n I
\end{equation}
donde \(n\) es el número de vueltas por unidad de longitud.

\section{Fuerza magnética entre corrientes}

\subsection{Fuerza entre dos conductores paralelos}
La fuerza por unidad de longitud entre dos conductores paralelos separados por una distancia \(d\) es:
\begin{equation}
    \frac{F}{L} = \frac{\mu_0 I_1 I_2}{2\pi d}
\end{equation}
\begin{itemize}
    \item Atracción: corrientes en el mismo sentido.
    \item Repulsión: corrientes en sentidos opuestos.
\end{itemize}

\resumebox{ejercicio}{Fuerza magnética del protón}{\footnotesize
    Un protón se mueve a \(2 \times 10^6 \, \text{m/s}\) perpendicular a un campo magnético de \(0.5 \, \text{T}\). Calcula la fuerza magnética sobre él.
}